%=========================================================
% Notas para el curso de Calculo Integral - TNM
% Formato conservado de Unidad2_Calculo_Integral_TNM
% Referencia directa: Aplicaciones_Integracion.pdf, cap. 6 de Stewart (6.1--6.5).
% La motivacion formal y las pruebas siguen una exposicion rigurosa tipo Leithold.
%=========================================================
\documentclass[12pt]{article}
\usepackage[utf8]{inputenc}
\usepackage[T1]{fontenc}
\usepackage[spanish,mexico]{babel}
\usepackage{lmodern}
\usepackage{amsmath,amssymb,amsthm,mathtools}
\usepackage{tikz}
\usepackage{pgfplots}
\pgfplotsset{compat=1.18}
\usepgfplotslibrary{fillbetween}
\usepackage[most]{tcolorbox}
\usepackage{multicol}
\usepackage{needspace}
\usepackage{geometry,enumitem,booktabs,array,xcolor,graphicx,hyperref,fancyhdr,microtype,longtable}
\geometry{letterpaper,top=2.3cm,bottom=2.3cm,left=2.5cm,right=2.5cm}
\setlength{\parindent}{0pt}
\setlength{\parskip}{0.55em}
\setlength{\headheight}{15pt}
\setlength{\emergencystretch}{2em}
\definecolor{azul}{RGB}{25,70,120}
\hypersetup{colorlinks=true,linkcolor=azul,urlcolor=azul,citecolor=azul,pdftitle={Calculo Integral - Unidad 3: Aplicaciones de la integral},pdfauthor={Carlos Ernesto Martinez Rodriguez}}
\newtheorem{teorema}{Teorema}[section]
\newtheorem{proposicion}[teorema]{Proposición}
\newtheorem{corolario}[teorema]{Corolario}
\theoremstyle{definition}
\newtheorem{definicion}[teorema]{Definición}
\newtheorem{ejemplo}[teorema]{Ejemplo}
\newtheorem{Ejer}[teorema]{Ejercicio}
\newtheorem{nota}[teorema]{Nota}
\theoremstyle{remark}
\newtheorem*{observacion}{Observación}
\newtheorem*{advertencia}{Advertencia}
\newtheorem*{estrategia}{Estrategia}
\newtcolorbox{notateorica}[1][]{enhanced,breakable,colback=gray!5,colframe=black!65,fonttitle=\bfseries,title=#1,boxrule=0.7pt,arc=2mm,left=2mm,right=2mm,top=2mm,bottom=2mm}
\newtcolorbox{formula}[1][]{enhanced,breakable,colback=blue!3,colframe=blue!55!black,fonttitle=\bfseries,title=#1,boxrule=0.8pt,arc=2mm,left=2mm,right=2mm,top=2mm,bottom=2mm}
\newcommand{\R}{\mathbb{R}}
\newcommand{\dd}{\,dx}
\newcommand{\dt}{\,dt}
\newcommand{\du}{\,du}
\newcommand{\diff}{\mathop{}\!\mathrm{d}}
\newcommand{\F}[1]{\[\displaystyle #1\]}
\pagestyle{fancy}
\fancyhf{}
\fancyhead[L]{\small Tecnológico Nacional de México}
\fancyhead[R]{\small Cálculo Integral}
\fancyfoot[L]{\footnotesize Carlos Ernesto Martínez Rodríguez}
\fancyfoot[C]{\footnotesize\thepage}
\fancyfoot[R]{\footnotesize mrdgzcarlose@gmail.com}
\renewcommand{\headrulewidth}{0.4pt}
\renewcommand{\footrulewidth}{0.4pt}
\title{\textbf{Notas para el curso de Cálculo Integral}\\[0.3cm]\Large Unidad 3: Aplicaciones de la integral}
\author{Carlos Ernesto Martínez Rodríguez\\Tecnológico Nacional de México}
\date{Agosto--Diciembre 2026\\Versión ampliada: noviembre de 2026}
\begin{document}
\maketitle
\tableofcontents
\newpage
%=========================================================
\section*{Presentación}
%=========================================================
En la Unidad 1 construimos la integral definida como límite de sumas de Riemann y demostramos el Teorema Fundamental del Cálculo. En la Unidad 2 desarrollamos técnicas de integración. Ahora utilizaremos esas herramientas para \emph{plantear} y resolver problemas: el primer paso no será integrar, sino determinar qué se está acumulando y respecto de qué variable. Distinguiremos claramente entre una fórmula deducida y una receta de cálculo, y justificaremos los signos, las dimensiones y la elección de cortes en cada problema.

Seguiremos una metodología de construcción geométrica, central en el enfoque de Stewart, y un desarrollo de las justificaciones con el detalle algebraico que caracteriza a Leithold: identificar una cantidad elemental, escribir una suma aproximante y pasar al límite. Siempre que sea posible, una misma región o sólido será estudiado por dos métodos, para comprender por qué aparecen fórmulas diferentes y cuándo conviene cada una.

\textbf{Propósito del temario:} calcular áreas, volúmenes de revolución, longitudes, superficies, trabajo, centros de masa e integrales impropias; interpretar y verificar aplicaciones en ingeniería. La planeación prevista para esta unidad es de 20 horas (12 teóricas y 8 prácticas).

%=========================================================

\subsection*{Un procedimiento general para construir modelos de integración}
En todos los temas partiremos de la misma idea: dividir una cantidad física o geométrica en partes pequeñas, aproximar cada parte mediante una expresión sencilla y sumar. Si la partición se refina, el límite de la suma define una integral:
\[
Q\approx\sum_{i=1}^n q(x_i^*)\,\Delta x_i,
\qquad Q=\lim_{\|P\|\to0}\sum_{i=1}^n q(x_i^*)\,\Delta x_i
=\int_a^b q(x)\,dx.
\]
La expresión de $q$ no debe memorizarse aisladamente: se deduce del \emph{elemento representativo}. En cada problema identificaremos (i) la región o sistema, (ii) la variable y el intervalo, (iii) el tamaño y significado de una porción pequeña, (iv) la integral, (v) las unidades y una comprobación geométrica o física.

\begin{ejemplo}[Del área elemental a la integral: un primer modelo]
Sea una región bajo $y=3-x$, con $0\leq x\leq2$. Una tira vertical tiene aproximadamente ancho $\Delta x$ y altura $3-x^*$; entonces $\Delta A\approx (3-x^*)\Delta x$. Al sumar las contribuciones y hacer fina la partición,
\[
A=\int_0^2(3-x)\,dx
=\left[3x-\frac{x^2}{2}\right]_0^2=4.
\]
El resultado coincide con el área del trapecio de bases $3$ y $1$ y anchura $2$: $(3+1)2/2=4$. La comprobación elemental confirma la construcción.
\end{ejemplo}

\begin{observacion}
La \emph{variable de integración} no es un detalle formal: indica cómo se ha cortado la región. En problemas de volumen se buscan unidades cúbicas; en superficies, unidades cuadradas; en trabajo, fuerza por distancia. Si las dimensiones no coinciden, el planteamiento debe revisarse antes de integrar.
\end{observacion}

%---------------------------------------------------------


\section{Área bajo una curva y área entre curvas}
%=========================================================
\subsection*{De la acumulación con signo al área geométrica}
Sea $f:[a,b]\to\R$ continua. La integral $\int_a^b f(x)\dd$ registra acumulación con signo: los sectores bajo el eje $x$ contribuyen negativamente. El \emph{área geométrica} nunca es negativa.
\begin{proposicion}
Si $f$ es continua, el área entre $y=f(x)$ y el eje $x$ sobre $[a,b]$ es
\[\boxed{A=\int_a^b |f(x)|\dd.}\]
Si $f\geq0$ sobre el intervalo, se simplifica a $A=\int_a^b f(x)\dd$.
\end{proposicion}
La razón es la misma que en la construcción de Riemann: un rectángulo de anchura $\Delta x$ y altura geométrica $|f(x_i^*)|$ tiene área $|f(x_i^*)|\Delta x$. Al sumar y refinar la partición aparece la integral del valor absoluto.

\begin{ejemplo}[Cuando una integral no es un área]
Calcular el área entre $y=x-1$ y el eje $x$ para $0\leq x\leq3$. La función cambia de signo en $x=1$, de manera que el área debe dividirse:
\begin{align*}
A&=\int_0^1(1-x)\dd+\int_1^3(x-1)\dd\\
 &=\left[x-\frac{x^2}{2}\right]_0^1+\left[\frac{x^2}{2}-x\right]_1^3
 =\frac12+2=\boxed{\frac52}.
\end{align*}
En cambio, $\int_0^3(x-1)\dd=3/2$ mide acumulación \emph{neta}, no área geométrica.
\end{ejemplo}

\subsection*{Área entre dos gráficas: rebanadas verticales}
Supongamos que en $[a,b]$ las funciones continuas $f$ y $g$ satisfacen $f(x)\geq g(x)$. Al fijar una abscisa $x$, una tira vertical de espesor $\Delta x$ tiene altura aproximada $f(x)-g(x)$; así,
\[A\approx\sum_{i=1}^n\bigl[f(x_i^*)-g(x_i^*)\bigr]\Delta x.\]
Al tomar el límite obtenemos la regla fundamental:
\begin{formula}[Área mediante rebanadas verticales]
\[\boxed{A=\int_a^b\bigl[f(x)-g(x)\bigr]\dd
=\int_a^b(\text{superior}-\text{inferior})\dd.}\]
Si las gráficas se cruzan, se localizan primero los puntos de intersección y se divide el intervalo cuando cambia cuál función queda arriba.
\end{formula}

\begin{ejemplo}[Una región limitada por dos parábolas]
Calcular el área encerrada por $y=x$ y $y=x^2$. Las intersecciones satisfacen $x=x^2$, es decir $x=0$ o $x=1$. En $0<x<1$ se cumple $x>x^2$. Por tanto
\[A=\int_0^1(x-x^2)\dd
=\left[\frac{x^2}{2}-\frac{x^3}{3}\right]_0^1
=\boxed{\frac16}.\]
La integral es positiva porque restamos la altura inferior a la superior.
\end{ejemplo}
\begin{center}
\begin{tikzpicture}
\begin{axis}[width=.70\linewidth,height=5.4cm,axis lines=middle,xmin=-.12,xmax=1.22,ymin=-.08,ymax=1.27,xtick={0,1},ytick={0,1},xlabel={$x$},ylabel={$y$},legend style={at={(.03,.97)},anchor=north west,font=\small,draw=none}]
\addplot[name path=recta,domain=0:1.1,samples=80,very thick,blue] {x};
\addlegendentry{$y=x$}
\addplot[name path=parabola,domain=0:1.1,samples=80,thick,red!75!black] {x^2};
\addlegendentry{$y=x^2$}
\addplot[blue!15] fill between[of=recta and parabola,soft clip={domain=0:1}];
\end{axis}
\end{tikzpicture}

{\small Región entre $y=x$ (arriba) y $y=x^2$ (abajo) para $0\leq x\leq1$.}
\end{center}

\begin{ejemplo}[El papel de los puntos de cruce]
Calcular el área entre $f(x)=x$ y $g(x)=x^3$ para $-1\leq x\leq1$. Las curvas se cortan en $-1,0,1$. En $(-1,0)$ se tiene $x^3>x$, mientras que en $(0,1)$ ocurre $x>x^3$. Así,
\begin{align*}
A&=\int_{-1}^0(x^3-x)\dd+\int_0^1(x-x^3)\dd\\
&=\left[\frac{x^4}{4}-\frac{x^2}{2}\right]_{-1}^0
+\left[\frac{x^2}{2}-\frac{x^4}{4}\right]_0^1
=\frac14+\frac14=\boxed{\frac12}.
\end{align*}
Una única integral de $x-x^3$ entre $-1$ y $1$ daría cero por cancelación.
\end{ejemplo}


\subsection*{Cuando las curvas se intersectan: localizar dónde cambia el orden}
La región entre dos funciones no siempre puede calcularse con una única resta. Si cambia la curva superior, debe partirse el intervalo en cada intersección pertinente. De manera compacta, para $f,g$ continuas,
\[
A=\int_a^b |f(x)-g(x)|\,dx.
\]
En la práctica conviene quitar el valor absoluto por intervalos: la integral del valor absoluto es una \emph{herramienta de comprobación}, no un sustituto del análisis geométrico.
\begin{ejemplo}[Cruce de funciones trigonométricas]
Determinar el área comprendida entre $y=\sin x$ y $y=\cos x$ para $0\le x\le\pi/2$. Primero resolvemos $\sin x=\cos x$, que en el intervalo dado implica $x=\pi/4$. Para $x<\pi/4$ la curva superior es $\cos x$; después del cruce es $\sin x$. Por tanto,
\begin{align*}
A&=\int_0^{\pi/4}(\cos x-\sin x)\,dx
+\int_{\pi/4}^{\pi/2}(\sin x-\cos x)\,dx\\
&=[\sin x+\cos x]_0^{\pi/4}
+[-\cos x-\sin x]_{\pi/4}^{\pi/2}\\
&=(\sqrt2-1)+(\sqrt2-1)
=\boxed{2\sqrt2-2}.
\end{align*}
Una sola resta $\int_0^{\pi/2}(\sin x-\cos x)\,dx$ daría $0$ por cancelación, no por ausencia de área.
\end{ejemplo}
\begin{proposicion}[Área nula entre curvas]
Sean $a<b$ y $f,g$ continuas en $[a,b]$. Si $A=\int_a^b|f-g|\,dx=0$, entonces $f(x)=g(x)$ para toda $x\in[a,b]$.
\end{proposicion}
\begin{proof}
Si en algún $c$ ocurriera $|f(c)-g(c)|>0$, por continuidad existiría un pequeño intervalo alrededor de $c$ en que $|f-g|$ sería mayor que una constante positiva. La integral en ese intervalo sería positiva, contradiciendo que el área total sea cero.
\end{proof}


\subsection*{Rebanadas horizontales e integración respecto de $y$}
Una región puede describirse mejor por funciones $x=r(y)$ y $x=\ell(y)$ para $c\le y\le d$, con $r(y)\ge\ell(y)$. Entonces cada tira horizontal tiene longitud $r(y)-\ell(y)$ y espesor $\Delta y$.
\begin{formula}[Área mediante rebanadas horizontales]
\[\boxed{A=\int_c^d\bigl[r(y)-\ell(y)\bigr]\,dy
=\int_c^d(\text{derecha}-\text{izquierda})\,dy.}\]
\end{formula}
\begin{ejemplo}[Elegir la variable más conveniente]
La región está limitada por $x=y^2$ y $x=2y$ en el primer cuadrante. Sus cruces cumplen $y^2=2y$, luego $y=0$ y $y=2$. En ese intervalo la frontera derecha es $x=2y$ y la izquierda $x=y^2$. Entonces
\[A=\int_0^2(2y-y^2)\,dy
=\left[y^2-\frac{y^3}{3}\right]_0^2
=\boxed{\frac43}.\]
Despejar $y$ para integrar respecto de $x$ también es posible, pero exigiría tener en cuenta ramas de la raíz. La elección de $dy$ simplifica el planteamiento.
\end{ejemplo}

\begin{estrategia}
Antes de integrar un área: (1) bosqueje las curvas; (2) calcule intersecciones; (3) elija tiras verticales u horizontales; (4) identifique arriba/abajo o derecha/izquierda; (5) divida si cambia el orden; (6) verifique que el resultado sea no negativo.
\end{estrategia}

%=========================================================

\subsection*{Áreas aproximadas a partir de mediciones: regla del punto medio}
En una aplicación de ingeniería, las fronteras o las velocidades pueden conocerse sólo mediante datos. En tal caso se puede \emph{estimar} la integral por una suma de rectángulos de punto medio. Para $n$ subintervalos iguales, $\Delta x=(b-a)/n$ y $m_i=a+(i-\tfrac12)\Delta x$,
\[
M_n=\Delta x\sum_{i=1}^n f(m_i)\approx\int_a^b f(x)\,dx.
\]
La misma idea aproxima áreas entre curvas si $f(m_i)$ se sustituye por la separación vertical positiva.
\begin{ejemplo}[Interpretación física de un área entre gráficas de velocidad]
Dos vehículos parten del mismo punto y avanzan en el mismo sentido. Se registran las diferencias $v_A-v_B$ de sus velocidades (en $\mathrm{m/s}$) en los puntos medios de cuatro intervalos de duración $2\,\mathrm s$:
\[
\begin{array}{c|cccc}
t\ (\mathrm s) &1&3&5&7\\ \hline
v_A-v_B\ (\mathrm{m/s}) &1&3&2&1
\end{array}
\]
Si $A$ es más rápido en todo $[0,8]$, la separación al cabo de $8$ segundos es aproximadamente
\[
\int_0^8(v_A-v_B)\,dt\approx2(1+3+2+1)=\boxed{14\,\mathrm m}.
\]
No se trata de una velocidad promedio sino de una \emph{distancia acumulada}. Al multiplicar $\mathrm{m/s}$ por segundos se obtienen metros.
\end{ejemplo}

%=========================================================


\section{Volúmenes por secciones y sólidos de revolución}
%=========================================================
\subsection*{Por qué integrar áreas transversales}
Imaginemos un sólido cortado por planos perpendiculares al eje $x$. Si la sección a la altura $x$ tiene área $S(x)$, una lámina delgada de espesor $\Delta x$ tiene volumen aproximado $S(x)\Delta x$. La suma y el límite de Riemann conducen a
\begin{teorema}[Volumen por secciones transversales]
Si $S$ es continua en $[a,b]$,
\[\boxed{V=\int_a^b S(x)\dd.}\]
La misma expresión es válida respecto de $y$ cuando las secciones son horizontales.
\end{teorema}
No se debe comenzar con $\pi\int f(x)^2\dd$ sin justificar primero la forma de la sección: esa fórmula sólo corresponde a discos de radio $f(x)$ respecto del eje de giro elegido.


\subsection*{Secciones transversales que no son circulares}
La fórmula $V=\int A(x)\,dx$ no exige que el sólido sea de revolución. Basta conocer cómo varía el área $A(x)$ de sus cortes perpendiculares al eje $x$. En ejercicios de este tipo conviene \emph{dibujar la base}, identificar en ella el segmento que determina el tamaño de cada sección y utilizar la fórmula geométrica de esa sección.

\begin{center}
\begin{tikzpicture}[scale=1.05]
\fill[blue!9] (0,0) -- plot[domain=0:4,samples=50] (\x,{sqrt(\x)}) -- (4,0) -- cycle;
\draw[->] (-.2,0)--(4.7,0) node[right] {$x$};
\draw[->] (0,-.1)--(0,2.6) node[above] {$y$};
\draw[blue!75!black,very thick,domain=0:4,samples=40] plot (\x,{sqrt(\x)});
\draw[red!75!black,very thick] (2.25,0)--(2.25,1.5);
\fill[red!75!black] (2.25,1.5) circle(1.8pt);
\node[anchor=west] at (2.35,1.5) {$s(x)=\sqrt{x}$};
\node[below] at (2.25,0) {$x$};
\node[below] at (4,0) {$4$};
\node at (1.0,.45) {Base};
\end{tikzpicture}

{\small La cuerda vertical de la base determina el lado de la sección cuadrada.}
\end{center}

\begin{ejemplo}[Secciones cuadradas: construcción completa]
La base de un sólido es la región limitada por $y=\sqrt x$, el eje $x$, $x=0$ y $x=4$. Las secciones perpendiculares al eje $x$ son cuadrados. Para un corte en $x$, el lado del cuadrado es la distancia vertical desde $0$ hasta $\sqrt x$, es decir, $s(x)=\sqrt x$. En consecuencia,
\[
A(x)=s(x)^2=x,
\qquad
V=\int_0^4x\,dx=\left[\frac{x^2}{2}\right]_0^4
=\boxed{8}.
\]
Si las coordenadas se miden en centímetros, el volumen se expresa en $\mathrm{cm^3}$.
\end{ejemplo}
\begin{ejemplo}[Secciones triangulares equiláteras]
La base es la región bajo $y=x$ en $0\le x\le2$. Cada sección perpendicular al eje $x$ es un triángulo equilátero cuyo lado mide $s(x)=x$. Recordando que el área de un triángulo equilátero es $\sqrt3\,s^2/4$, tenemos
\[
A(x)=\frac{\sqrt3}{4}x^2,
\qquad
V=\frac{\sqrt3}{4}\int_0^2x^2\,dx
=\boxed{\frac{2\sqrt3}{3}}.
\]
Aquí no aparece $\pi$: las secciones no son discos.
\end{ejemplo}
\begin{ejemplo}[Secciones semicirculares]
La base es el disco $x^2+y^2\le4$. Cada corte perpendicular al eje $x$ es un semicírculo cuyo \emph{diámetro} es la cuerda vertical de la base: $d(x)=2\sqrt{4-x^2}$. El radio de la sección es, por tanto, $r(x)=\sqrt{4-x^2}$, y
\[
A(x)=\frac\pi2(4-x^2).
\]
Al integrar entre $-2$ y $2$,
\[
V=\frac\pi2\int_{-2}^2(4-x^2)\,dx
=\frac\pi2\left[4x-\frac{x^3}{3}\right]_{-2}^2
=\boxed{\frac{16\pi}{3}}.
\]
Una confusión frecuente consiste en usar el diámetro como si fuera el radio.
\end{ejemplo}
\begin{estrategia}
Para secciones conocidas, ordene el procedimiento como sigue: \emph{base plana} $\longrightarrow$ \emph{longitud de la sección} $\longrightarrow$ \emph{área transversal} $A(x)$ $\longrightarrow$ \emph{integral del volumen}. No suponga que todos los problemas son sólidos de revolución.
\end{estrategia}



\subsection*{Método de discos}
Al girar la región comprendida entre $y=f(x)\geq0$ y el eje $x$ alrededor de este último, cada sección perpendicular al eje es un disco de radio $R(x)=f(x)$; por tanto $S(x)=\pi R(x)^2$.
\begin{formula}[Discos]
\[\boxed{V=\pi\int_a^b [R(x)]^2\dd.}\]
Para un giro alrededor del eje $y$ con cortes perpendiculares al eje de giro se integra respecto de $y$.
\end{formula}
\begin{ejemplo}[Rotación de una parábola]
La región bajo $y=x^2$ para $0\leq x\leq2$ gira alrededor del eje $x$. El radio exterior en la abscisa $x$ es $R(x)=x^2$ y no existe hueco central. Por eso
\[V=\pi\int_0^2(x^2)^2\dd
=\pi\left[\frac{x^5}{5}\right]_0^2=\boxed{\frac{32\pi}{5}}.\]
La unidad resultante es cúbica, a diferencia de la integral utilizada para área.
\end{ejemplo}
\begin{ejemplo}[Discos alrededor del eje $y$]
Considere la región del primer cuadrante comprendida entre $x=\sqrt y$, el eje $y$ y $0\le y\le4$. Gira alrededor del eje $y$. Una sección horizontal es un disco de radio $R(y)=\sqrt y$:
\[V=\pi\int_0^4(\sqrt y)^2\,dy
=\pi\left[\frac{y^2}{2}\right]_0^4=\boxed{8\pi}.\]
\end{ejemplo}

\subsection*{Método de arandelas}
Cuando la región no toca el eje de giro, la sección transversal tiene un agujero: es una arandela. Su área es la diferencia entre las áreas de dos discos, \emph{no} el cuadrado de la diferencia de radios.
\begin{formula}[Arandelas]
\[\boxed{V=\pi\int_a^b\bigl(R(x)^2-r(x)^2\bigr)\dd,\qquad R(x)\ge r(x)\ge0.}\]
Los radios son distancias al eje de rotación. Si el eje es $y=k$, deben medirse como $|f(x)-k|$ y $|g(x)-k|$, escogiendo correctamente el mayor y el menor.
\end{formula}
\begin{ejemplo}[Arandelas alrededor del eje $x$]
La región entre $y=2$ y $y=x$ para $0\leq x\leq2$ gira alrededor del eje $x$. Para cada $x$, el radio exterior es $R(x)=2$ y el interior $r(x)=x$. Por tanto,
\begin{align*}
V&=\pi\int_0^2(2^2-x^2)\dd
=\pi\left[4x-\frac{x^3}{3}\right]_0^2\\
&=\boxed{\frac{16\pi}{3}}.
\end{align*}
La expresión incorrecta $\pi\int_0^2(2-x)^2\dd$ confundiría el grosor de la arandela con su área.
\end{ejemplo}
\begin{ejemplo}[El eje de giro está desplazado]
Considere la región entre $y=x$ y $y=x^2$ para $0\le x\le1$, girada alrededor de $y=-1$. Ambos radios se miden \emph{desde} esa recta: $R(x)=x+1$ y $r(x)=x^2+1$. De aquí,
\begin{align*}
V&=\pi\int_0^1\bigl[(x+1)^2-(x^2+1)^2\bigr]\dd\\
&=\pi\int_0^1(-x^4-x^2+2x)\dd
=\pi\left(-\frac15-\frac13+1\right)
=\boxed{\frac{7\pi}{15}}.
\end{align*}
\end{ejemplo}
\begin{observacion}
Si la región cruza el eje de giro, los radios requieren una inspección geométrica adicional: las secciones pueden ser discos, arandelas o cambiar de tipo. El procedimiento de radios ``superior menos inferior'' no es universal cuando el eje atraviesa la región.
\end{observacion}


\begin{center}
\begin{tikzpicture}[scale=2.7]
\fill[blue!13] plot[domain=0:1,samples=40] (\x,\x) -- plot[domain=1:0,samples=40] (\x,{\x*\x}) -- cycle;
\draw[->] (-.07,0)--(1.24,0) node[right] {$x$};
\draw[->] (0,-.04)--(0,2.18) node[above] {$y$};
\draw[blue!70!black,thick,domain=0:1] plot (\x,\x);
\draw[red!70!black,thick,domain=0:1] plot (\x,{\x*\x});
\draw[dashed,thick] (-.08,2)--(1.2,2) node[right] {$y=2$};
\draw[<->,red!70!black] (.45,.2025)--(.45,2);
\draw[<->,blue!70!black] (.68,.68)--(.68,2);
\node[anchor=east,fill=white,inner sep=1pt,scale=.62] at (.40,1.42) {$R=2-x^2$};
\node[anchor=west,fill=white,inner sep=1pt,scale=.62] at (.70,1.25) {$r=2-x$};
\end{tikzpicture}

{\small El radio exterior llega a la curva inferior al girar respecto a la recta $y=2$.}
\end{center}

\begin{ejemplo}[Arandelas cuando el eje de giro está arriba de la región]
Giremos alrededor de $y=2$ la región entre $y=x$ y $y=x^2$, con $0\le x\le1$. Como el eje está por encima de ambas curvas, el radio \emph{exterior} corresponde a la curva \emph{inferior} $y=x^2$:
\[
R(x)=2-x^2,\qquad r(x)=2-x.
\]
Por ello,
\begin{align*}
V&=\pi\int_0^1\bigl[(2-x^2)^2-(2-x)^2\bigr]dx\\
&=\pi\int_0^1\bigl(x^4-5x^2+4x\bigr)dx\\
&=\pi\left[\frac{x^5}{5}-\frac{5x^3}{3}+2x^2\right]_0^1
=\boxed{\frac{8\pi}{15}}.
\end{align*}
Al revisar el dibujo debe ser claro por qué la distancia a $y=2$ invierte el orden de los radios.
\end{ejemplo}
\begin{ejemplo}[Una región que atraviesa el eje de giro]
El rectángulo $0\le x\le2$, $-1\le y\le1$ gira alrededor del eje $x$. En cada abscisa su sección transversal es un \emph{disco lleno} de radio $1$, no una arandela con radios $1$ y $1$. Por tanto,
\[V=\pi\int_0^2 1^2\,dx=\boxed{2\pi}.\]
Restar mecánicamente los cuadrados de las ordenadas superior e inferior produciría el resultado falso $0$. Las distancias se determinan sobre la sección \emph{real del sólido}.
\end{ejemplo}



\subsection*{Método de cascarones cilíndricos}
En lugar de cortar perpendicularmente al eje de giro, tomamos una tira \emph{paralela} a ese eje. Al girarla se forma un cascarón cilíndrico de radio $r$, altura $h$ y espesor $\Delta r$. La diferencia de volúmenes de cilindros es
\[\pi(r+\Delta r)^2h-\pi r^2h=2\pi rh\Delta r+\pi h(\Delta r)^2.\]
En el límite, el término cuadrático en el espesor deja de contribuir y obtenemos
\begin{formula}[Cascarones cilíndricos]
\[\boxed{V=2\pi\int_a^b (\text{radio})(\text{altura})\dd.}\]
Para girar alrededor del eje $y$ usando tiras verticales, $r=x$ y $h=f(x)-g(x)$ cuando la región está a la derecha del eje. Para girar alrededor del eje $x$ usando tiras horizontales, la integración es en $dy$.
\end{formula}
\begin{center}
\begin{tikzpicture}[scale=1.22]
\fill[blue!9] (0,0) -- plot[domain=0:2,samples=50] (\x,{0.45*\x*\x}) -- (2,0) -- cycle;
\draw[->] (-.18,0) -- (2.6,0) node[right] {$x$};
\draw[->] (0,-.12) -- (0,2.2) node[above] {$y$};
\draw[blue!75!black,very thick] plot[domain=0:2,samples=45] (\x,{0.45*\x*\x});
\fill[red!30] (1.32,0) rectangle (1.44,.87);
\draw[red!70!black,thick] (1.32,0) rectangle (1.44,.87);
\draw[<->] (.02,-.22) -- (1.38,-.22) node[midway,below] {$r=x$};
\draw[<->] (1.60,0) -- (1.60,.87) node[midway,right] {$h=f(x)$};
\node[anchor=west,align=left] at (2.60,1.48) {\small Tira paralela\\\small al eje $y$\\$\Downarrow$ cascarón};
\draw[->,red!70!black] (2.58,1.08) -- (1.42,.62);
\end{tikzpicture}

{\small Esquema de la tira vertical utilizada al girar una región alrededor del eje $y$.}
\end{center}

\begin{ejemplo}[Cascarones: región bajo una parábola]
Gire alrededor del eje $y$ la región delimitada por $y=x^2$, $y=0$ y $x=2$. Una tira vertical en $x$ recorre una altura $h=x^2$ y describe un cascarón de radio $r=x$:
\[V=2\pi\int_0^2 x(x^2)\dd
=2\pi\left[\frac{x^4}{4}\right]_0^2
=\boxed{8\pi}.\]
Si usáramos arandelas horizontales habría que describir la región de otra forma: para $0\le y\le4$, el radio exterior es $2$ y el interior es $\sqrt y$. Así,
\[V=\pi\int_0^4(4-y)\,dy
=\pi[4y-y^2/2]_0^4=8\pi.\]
La coincidencia de resultados valida ambos planteamientos.
\end{ejemplo}
\begin{ejemplo}[Cascarones alrededor de una recta vertical]
La región bajo $y=x$ entre $x=0$ y $x=2$ gira alrededor de $x=3$. El radio de una tira es $r(x)=3-x$, y su altura $h(x)=x$. Por tanto
\[V=2\pi\int_0^2(3-x)x\dd
=2\pi\left[\frac{3x^2}{2}-\frac{x^3}{3}\right]_0^2
=\boxed{\frac{20\pi}{3}}.\]
La distancia al eje no es $x$: el eje de rotación se encuentra en $x=3$.
\end{ejemplo}
\begin{estrategia}
\textbf{Elección del método.} Para cada dibujo trace una tira de integración. Si la tira es perpendicular al eje de giro, use discos o arandelas. Si es paralela, use cascarones. La forma de la tira determina la fórmula, no una preferencia previa por $dx$ o $dy$.
\end{estrategia}

%=========================================================

\begin{ejemplo}[Cascarones horizontales y verificación por arandelas]
La región delimitada por $x=y^2$ y $x=y$ para $0\le y\le1$ gira alrededor del eje $x$. Una tira horizontal tiene altura (longitud horizontal) $h(y)=y-y^2$, radio $r(y)=y$ y espesor $dy$. El método de cascarones da
\[
V=2\pi\int_0^1 y(y-y^2)\,dy
=2\pi\left(\frac13-\frac14\right)
=\boxed{\frac\pi6}.
\]
Para verificarlo, describamos la misma región mediante cortes verticales: para $0\le x\le1$ se cumple $x\le y\le\sqrt x$. Los radios son $R(x)=\sqrt x$ y $r(x)=x$, de donde
\[
V=\pi\int_0^1\bigl[(\sqrt x)^2-x^2\bigr]dx
=\pi\left(\frac12-\frac13\right)
=\frac\pi6.
\]
Los dos métodos entregan el mismo volumen porque describen el mismo conjunto tridimensional, aunque emplean variables de integración distintas.
\end{ejemplo}

%=========================================================


\section{Longitud de arco}
%=========================================================
Una curva no puede medirse sumando distancias horizontales: debe considerarse simultáneamente su desplazamiento en $x$ y en $y$. Aproximemos un arco de $y=f(x)$, con $f'$ continua en $[a,b]$, mediante segmentos rectos entre puntos consecutivos de una partición. Por Pitágoras, su longitud es
\[\Delta s_i=\sqrt{(\Delta x_i)^2+(\Delta y_i)^2}
=\sqrt{1+\left(\frac{\Delta y_i}{\Delta x_i}\right)^2}\,\Delta x_i.\]
El Teorema del Valor Medio garantiza la existencia de $c_i$ entre $x_{i-1}$ y $x_i$ tal que $\Delta y_i/\Delta x_i=f'(c_i)$. Las longitudes poligonales se aproximan a
\[\sum_i\sqrt{1+[f'(c_i)]^2}\,\Delta x_i,\]
una suma de Riemann.
\begin{teorema}[Longitud de arco]
Si $f'$ es continua sobre $[a,b]$, la longitud del arco de $y=f(x)$ entre $x=a$ y $x=b$ es
\[\boxed{L=\int_a^b\sqrt{1+[f'(x)]^2}\dd.}\]
Si se expresa $x=g(y)$ con $g'$ continua en $[c,d]$, entonces $L=\int_c^d\sqrt{1+[g'(y)]^2}\,dy$.
\end{teorema}
\begin{ejemplo}[Una recta como comprobación geométrica]
Encontrar la longitud de $y=2x+1$ sobre $0\le x\le3$. Como $f'(x)=2$,
\[L=\int_0^3\sqrt{1+4}\dd=3\sqrt5.\]
Por la fórmula de distancia, los extremos son $(0,1)$ y $(3,7)$, cuya distancia es $\sqrt{3^2+6^2}=3\sqrt5$. Ambos métodos coinciden.
\end{ejemplo}
\begin{ejemplo}[Una curva con integral elemental]
Calcule la longitud de $y=\frac23 x^{3/2}$ para $0\le x\le3$. La derivada es $f'(x)=\sqrt x$, de modo que
\[L=\int_0^3\sqrt{1+x}\dd
=\left[\frac23(1+x)^{3/2}\right]_0^3
=\boxed{\frac{14}{3}}.\]
La elección de esta función permitió que el integrando de longitud tuviera primitiva elemental; no todas las curvas poseen esa propiedad.
\end{ejemplo}
\begin{ejemplo}[Longitud de un cuarto de circunferencia]
La semicircunferencia superior de radio $r$ satisface $y=\sqrt{r^2-x^2}$. En $0\le x<r$,
\[y'=-\frac{x}{\sqrt{r^2-x^2}},\qquad
\sqrt{1+(y')^2}=\frac{r}{\sqrt{r^2-x^2}}.\]
Aunque la derivada crece sin cota al aproximarse $x$ a $r$, la longitud puede obtenerse mediante un límite:
\[L=\lim_{b\uparrow r}r\int_0^b\frac{dx}{\sqrt{r^2-x^2}}
=r\left[\arcsin\frac{x}{r}\right]_0^r
=\boxed{\frac{\pi r}{2}}.\]
La fórmula conocida de longitud de una circunferencia resulta así consistente con el cálculo integral.
\end{ejemplo}
\begin{observacion}
No debe confundirse $\int_a^b\sqrt{1+[f'(x)]^2}\dd$ con $\sqrt{1+\left[\int_a^b f'(x)\dd\right]^2}$. La raíz se aplica \emph{antes} de integrar: mide el incremento local de longitud.
\end{observacion}

%=========================================================

\subsection*{Cotas y longitud acumulada}
La deducción anterior, basada en el Teorema del Valor Medio y en sumas de Riemann, muestra cómo se construye la longitud. Podemos obtener además dos estimaciones útiles sin evaluar ninguna primitiva. Como $\sqrt{1+u^2}\ge1$ para todo $u$ real,
\[
L=\int_a^b\sqrt{1+[f'(x)]^2}\,dx\ge b-a.
\]
Si se conoce una cota $|f'(x)|\le M$ sobre $[a,b]$, entonces
\[
\boxed{b-a\le L\le(b-a)\sqrt{1+M^2}.}
\]
La primera desigualdad expresa que un recorrido sobre la curva no es más corto que su proyección horizontal; la segunda sirve para comprobar la razonabilidad del resultado.

Para medir sólo el arco desde el extremo izquierdo hasta una abscisa variable, definimos
\[
s(x)=\int_a^x\sqrt{1+[f'(t)]^2}\,dt.
\]
Por el Primer Teorema Fundamental del Cálculo, cuando $f'$ es continua,
\[
\boxed{s'(x)=\sqrt{1+[f'(x)]^2}.}
\]
Esta derivada representa la razón de cambio de la longitud acumulada con respecto a $x$.

\begin{ejemplo}[Una curva diseñada para simplificar el radical]
Calcule la longitud de la gráfica de
\[
y=\frac{x^2}{2}-\frac14\ln x,
\qquad 1\le x\le2.
\]
Primero derivamos:
\[
y'=x-\frac{1}{4x}.
\]
Al formar el radical observamos una identidad algebraica útil:
\[
1+(y')^2
=1+\left(x-\frac1{4x}\right)^2
=\left(x+\frac1{4x}\right)^2.
\]
Como $x>0$, la raíz positiva es $x+1/(4x)$ y, por tanto,
\begin{align*}
L&=\int_1^2\left(x+\frac1{4x}\right)dx
=\left[\frac{x^2}{2}+\frac14\ln x\right]_1^2\\
&=\boxed{\frac32+\frac14\ln2}.
\end{align*}
El ejemplo muestra por qué se debe \emph{simplificar el integrando antes de elegir una técnica de integración}.
\end{ejemplo}

%=========================================================


\section{Área de una superficie de revolución}
%=========================================================
Un segmento recto de longitud $\Delta s$, al rotar alrededor del eje $x$, genera aproximadamente la superficie lateral de un tronco de cono. Su área es $2\pi\bar r\,\Delta s$, donde $\bar r$ es el radio medio. Al refinar la partición, $\Delta s\approx\sqrt{1+[f'(x)]^2}\Delta x$ y $r=f(x)$ para $f\geq0$.
\begin{teorema}[Superficie de revolución]
Para una curva suave $y=f(x)\ge0$ sobre $[a,b]$, el área de superficie generada al girar alrededor del eje $x$ es
\[\boxed{S=2\pi\int_a^b f(x)\sqrt{1+[f'(x)]^2}\dd.}\]
Si la curva gira alrededor del eje $y$, la distancia al eje es $x$ (para $x\ge0$):
\[\boxed{S=2\pi\int_a^b x\sqrt{1+[f'(x)]^2}\dd.}\]
Con un eje desplazado debe sustituirse el radio por la distancia correspondiente.
\end{teorema}
\begin{ejemplo}[Superficie lateral de un cono]
Al girar $y=x$ para $0\le x\le2$ alrededor del eje $x$ se forma un cono de radio y altura $2$. Puesto que $f'=1$,
\[S=2\pi\int_0^2 x\sqrt2\dd
=2\pi\sqrt2\left[\frac{x^2}{2}\right]_0^2
=\boxed{4\pi\sqrt2}.\]
Geométricamente, $S=\pi r\ell=\pi(2)(2\sqrt2)$, pues la generatriz mide $\ell=2\sqrt2$. Esto verifica la fórmula.
\end{ejemplo}
\begin{ejemplo}[Rotación de $y=\sqrt x$]
Girar $y=\sqrt x$, para $1\le x\le4$, alrededor del eje $x$. Aquí $f'(x)=1/(2\sqrt x)$, de modo que
\begin{align*}
S&=2\pi\int_1^4\sqrt x\sqrt{1+\frac{1}{4x}}\dd
=\pi\int_1^4\sqrt{4x+1}\dd\\
&=\frac{\pi}{6}\left[(4x+1)^{3/2}\right]_1^4
=\boxed{\frac{\pi}{6}(17^{3/2}-5^{3/2})}.
\end{align*}
\end{ejemplo}
\begin{advertencia}
$S$ mide \emph{área superficial} y $V$ mide \emph{volumen}. Aunque ambos se obtienen por integración, no son magnitudes intercambiables. En un problema físico compruebe siempre las unidades: $S$ tiene unidades cuadradas y $V$ cúbicas.
\end{advertencia}

%=========================================================

\subsection*{De pequeños troncos de cono a superficies de revolución}
Si un arco corto de longitud $\Delta s$ se gira alrededor del eje $x$, produce aproximadamente la superficie lateral de un tronco de cono: $\Delta S\approx2\pi r^*\Delta s$, donde $r^*$ es un radio representativo. Al sumar y pasar al límite aparece $S=2\pi\int r\,ds$. Ésta es la razón geométrica por la que la fórmula de superficie no usa el \emph{cuadrado} del radio, a diferencia del volumen por discos.
\begin{ejemplo}[Superficie de una semiesfera]
La semicircunferencia superior $y=\sqrt{4-x^2}$, $0\le x\le2$, gira alrededor del eje $x$. Produce una \emph{semiesfera} de radio $2$. Para $0\le x<2$,
\[
y'=-\frac{x}{\sqrt{4-x^2}},
\qquad
\sqrt{1+(y')^2}=\frac2{\sqrt{4-x^2}}.
\]
Entonces el integrando de superficie simplifica a
\[
2\pi y\sqrt{1+(y')^2}
=2\pi\sqrt{4-x^2}\frac2{\sqrt{4-x^2}}=4\pi.
\]
Aunque $y'$ no está acotada cerca de $x=2$, evaluamos el límite desde la izquierda:
\[
S=\lim_{b\uparrow2}\int_0^b4\pi\,dx
=\boxed{8\pi}.
\]
Coincide con la mitad del área superficial de una esfera de radio $2$.
\end{ejemplo}
\begin{observacion}
Una superficie generada por un \emph{arco} no incluye automáticamente las tapas planas de un sólido. Si se pide el área total de un recipiente o cuerpo cerrado, deben agregarse por separado las superficies que falten.
\end{observacion}

%=========================================================


\section{Trabajo realizado por una fuerza variable}
%=========================================================
Para una fuerza constante $F$ paralela al desplazamiento $d$, el trabajo es $W=Fd$. Cuando la fuerza varía con la posición, se parte el desplazamiento en intervalos cortos. En cada uno, el trabajo se aproxima por $F(x_i^*)\Delta x_i$. El límite de las sumas da:
\begin{definicion}[Trabajo unidimensional]
Para una fuerza continua $F(x)$ en la dirección del eje $x$, el trabajo al desplazarse de $x=a$ a $x=b$ es
\[\boxed{W=\int_a^b F(x)\dd.}\]
Una fuerza opuesta al desplazamiento produce trabajo negativo; si sólo interesa el esfuerzo realizado, se debe interpretar el signo según el contexto.
\end{definicion}
\begin{ejemplo}[Resorte y ley de Hooke]
La fuerza necesaria para estirar un resorte ideal una distancia $x$ desde su longitud natural es $F(x)=kx$, donde $k>0$ es la constante elástica. Por ello,
\[W=\int_0^d kx\dd=\boxed{\frac12kd^2}.\]
Si $k=50\,\mathrm{N/m}$ y $d=0.20\,\mathrm m$,
\[W=\tfrac12(50)(0.20)^2=\boxed{1\,\mathrm J}.\]
El trabajo no es $F(d)d$, porque la fuerza comienza en cero y aumenta gradualmente.
\end{ejemplo}
\begin{ejemplo}[Trabajo para levantar un cable]
Un cable uniforme de longitud $L=4\,\mathrm m$ tiene peso por unidad de longitud $w=10\,\mathrm{N/m}$ y cuelga verticalmente desde un borde. Para subirlo completamente, una sección situada a distancia $y$ bajo el borde debe elevarse $y$ metros. Una tira de longitud $dy$ pesa $w\,dy$, por lo que
\[dW=w\,y\,dy.\]
Al sumar todas las tiras se obtiene
\[W=\int_0^4 10y\,dy=10\left[\frac{y^2}{2}\right]_0^4
=\boxed{80\,\mathrm J}.\]
\end{ejemplo}
\begin{ejemplo}[Bombeo de agua de un depósito]
Un depósito vertical tiene sección horizontal constante de $1\,\mathrm{m}^2$ y altura $2\,\mathrm m$; está lleno de agua. Tomemos $y$ como altura desde el fondo. Una lámina de espesor $dy$ tiene volumen $1\,\mathrm{m}^2\,dy$ y peso aproximado $1000(9.8)(1)\,dy$ newtons. Para llegar al borde superior debe subir $2-y$ metros. Entonces
\begin{align*}
W&=9800\int_0^2(2-y)\,dy\\
&=9800\left[2y-\frac{y^2}{2}\right]_0^2
=\boxed{19600\,\mathrm J}.
\end{align*}
Suponer que toda el agua sube $2$ metros duplicaría incorrectamente el resultado.
\end{ejemplo}

%=========================================================

\subsection*{El trabajo elemental como fuerza por desplazamiento}
En una dimensión, si la fuerza $F(x)$ actúa en el sentido positivo del desplazamiento, una porción pequeña realiza $\Delta W\approx F(x_i^*)\Delta x_i$ y, al sumar, $W=\int_a^b F(x)\,dx$. Si la fuerza se opone al movimiento, el trabajo \emph{de esa fuerza} es negativo. La fuerza externa que mantiene una velocidad constante realiza el trabajo positivo opuesto.
\begin{ejemplo}[Resistencia lineal al movimiento]
Un móvil se desplaza de $x=0$ a $x=3\,\mathrm m$ contra una fuerza de resistencia $F_r(x)=-4x\,\mathrm N$. El trabajo realizado por la resistencia es
\[
W_r=\int_0^3(-4x)\,dx=\boxed{-18\,\mathrm J}.
\]
Si no hay otras fuerzas que realicen trabajo neto y se mantiene velocidad constante, el agente externo debe suministrar $+18\,\mathrm J$. El signo forma parte del modelo físico.
\end{ejemplo}
\begin{ejemplo}[Tanque cuyo corte horizontal cambia con la altura]
Un tanque de profundidad $2\,\mathrm m$ contiene agua hasta el borde. A la altura $y$ medida desde el fondo, la sección horizontal tiene área $A(y)=2y\,\mathrm{m^2}$ ($0\le y\le2$). Cada lámina de volumen $A(y)\,dy$ tiene peso $\rho g A(y)\,dy$ y debe subir $2-y$ metros. Con $\rho=1000\,\mathrm{kg/m^3}$ y $g=9.8\,\mathrm{m/s^2}$,
\begin{align*}
W&=9800\int_0^2 2y(2-y)\,dy
=19600\left[y^2-\frac{y^3}{3}\right]_0^2\\
&=\boxed{\frac{78400}{3}\,\mathrm J}.
\end{align*}
Observe que el volumen de la lámina y la distancia de bombeo \emph{son funciones distintas de $y$}.
\end{ejemplo}
\subsection*{Una extensión de ingeniería: trabajo de un gas}
Si un gas ejerce presión $P(V)$ sobre un émbolo de área transversal $A$, el desplazamiento pequeño $dx$ realiza $dW=F\,dx=P(V)A\,dx=P(V)\,dV$. Así se obtiene
\[
\boxed{W=\int_{V_1}^{V_2}P(V)\,dV.}
\]
\begin{ejemplo}[Expansión a presión inversamente proporcional al volumen]
Suponga que en un proceso idealizado $P(V)=k/V$, con $k>0$, y el gas se expande desde $V_1$ hasta $V_2>V_1$. Entonces
\[
W=\int_{V_1}^{V_2}\frac{k}{V}\,dV
=\boxed{k\ln\!\left(\frac{V_2}{V_1}\right)}.
\]
En el sistema internacional la presión está en pascales y el volumen en $\mathrm{m^3}$; su producto es un joule.
\end{ejemplo}

%=========================================================
\section{Valor promedio de una función}
%=========================================================
Un conjunto finito de mediciones $y_1,\ldots,y_n$ tiene promedio $(y_1+\cdots+y_n)/n$. Si la magnitud varía continuamente, como ocurre con la temperatura durante el día o el consumo eléctrico, tomamos mediciones en subintervalos iguales de longitud $\Delta x=(b-a)/n$. El promedio de los valores muestreados puede reescribirse como
\[
\frac1n\sum_{i=1}^n f(x_i^*)
=\frac1{b-a}\sum_{i=1}^n f(x_i^*)\Delta x.
\]
El límite de estas sumas motiva la definición siguiente.
\begin{definicion}[Valor promedio en un intervalo]
Para $a<b$ y $f$ integrable en $[a,b]$, se define
\[
\boxed{f_{\mathrm{prom}}=\frac1{b-a}\int_a^b f(x)\,dx.}
\]
El valor promedio conserva las mismas unidades que $f$ porque la integral se divide entre la longitud del intervalo.
\end{definicion}
\begin{ejemplo}[Promedio de una función cuadrática]
Calcule el valor promedio de $f(x)=x^2$ en $[0,2]$. Aplicamos directamente la definición:
\[
f_{\mathrm{prom}}=\frac12\int_0^2x^2\,dx
=\frac12\left[\frac{x^3}{3}\right]_0^2
=\boxed{\frac43}.
\]
Geométricamente, el rectángulo de base $2$ y altura $4/3$ tiene la misma área que la región bajo $y=x^2$ entre $0$ y $2$. El promedio de los valores extremos sería $2$, un número \emph{diferente}.
\end{ejemplo}
\begin{teorema}[Valor medio para integrales]
Si $f$ es continua en $[a,b]$, con $a<b$, existe al menos un $c\in[a,b]$ tal que
\[
\boxed{f(c)=\frac1{b-a}\int_a^b f(x)\,dx.}
\]
\end{teorema}
\begin{proof}
Por continuidad, $f$ alcanza mínimo $m$ y máximo $M$ en $[a,b]$, luego $m\le f(x)\le M$. La monotonía de la integral produce
\[
m(b-a)\le\int_a^b f(x)\,dx\le M(b-a).
\]
Como $b-a>0$, al dividir se concluye $m\le f_{\mathrm{prom}}\le M$. Por el Teorema del Valor Intermedio, la función continua toma el valor $f_{\mathrm{prom}}$ en algún punto $c$.
\end{proof}
\begin{ejemplo}[Localización del instante que alcanza el promedio]
Para $f(x)=x^2$ en $[0,2]$, el promedio anterior es $4/3$. Buscamos $c$ tal que $f(c)=4/3$:
\[
c^2=\frac43,\qquad c=\boxed{\frac2{\sqrt3}}.
\]
Se toma la solución positiva porque $0\le c\le2$. El teorema garantiza al menos un punto y aquí es único.
\end{ejemplo}
\begin{ejemplo}[Velocidad promedio a partir de una función variable]
La velocidad de un vehículo durante $0\le t\le3\,\mathrm s$ se modela por $v(t)=2t+1$ $\mathrm{m/s}$. La distancia recorrida es
\[
D=\int_0^3(2t+1)\,dt=[t^2+t]_0^3=12\,\mathrm m.
\]
Por tanto, la velocidad promedio en ese periodo es
\[
v_{\mathrm{prom}}=\frac D3=\boxed{4\,\mathrm{m/s}}.
\]
No debemos confundir el valor medio de la velocidad con la media aritmética de lecturas tomadas en tiempos desiguales.
\end{ejemplo}

%=========================================================


\section{Centro de masa y momentos}
%=========================================================
La posición del centro de masa es un \emph{promedio ponderado}: las partes con más masa influyen más. Para partículas puntuales con masas $m_i$ en posiciones $x_i$, la coordenada horizontal es
\[\bar x=\frac{\sum_i m_i x_i}{\sum_i m_i}.\]
En una distribución continua las sumas se convierten en integrales.
\begin{definicion}[Varilla de densidad variable]
Una varilla delgada ocupa $a\le x\le b$ y posee densidad lineal $\rho(x)\ge0$. Su masa y primer momento respecto del origen son
\[m=\int_a^b\rho(x)\dd,\qquad M=\int_a^b x\rho(x)\dd.\]
Si $m>0$, su centro de masa se encuentra en
\[\boxed{\bar x=\frac{M}{m}.}\]
\end{definicion}
\begin{ejemplo}[Varilla no uniforme]
Una varilla ocupa $[0,2]$ y tiene densidad lineal $\rho(x)=1+x$ (unidades compatibles). Entonces
\[m=\int_0^2(1+x)\dd=\left[x+\frac{x^2}{2}\right]_0^2=4,\]
\[M=\int_0^2x(1+x)\dd
=\left[\frac{x^2}{2}+\frac{x^3}{3}\right]_0^2=\frac{14}{3}.\]
Su centro de masa es $\boxed{\bar x=\frac76}$. Está a la derecha del punto medio $x=1$ porque la densidad crece con $x$.
\end{ejemplo}

\subsection*{Láminas entre gráficas: momentos respecto de los ejes}
Una lámina homogénea de densidad superficial constante $\rho$ está situada entre $y=f(x)$ (arriba) y $y=g(x)$ (abajo), para $a\le x\le b$. Una tira vertical de anchura $dx$ tiene área $[f(x)-g(x)]dx$ y masa $dm=\rho[f-g]dx$. Su centro vertical está aproximadamente en $(f+g)/2$. Así,
\begin{formula}[Centroide de una región plana homogénea]
\begin{align*}
A&=\int_a^b[f(x)-g(x)]\dd,\\
M_y&=\rho\int_a^b x[f(x)-g(x)]\dd,\\
M_x&=\frac\rho2\int_a^b[f(x)^2-g(x)^2]\dd,\\
\bar x&=\frac{M_y}{\rho A},\qquad
\bar y=\frac{M_x}{\rho A}.
\end{align*}
Para densidad uniforme, el factor $\rho$ se cancela al calcular las coordenadas del centroide.
\end{formula}
\begin{ejemplo}[Centroide bajo la recta $y=x$]
La región triangular comprendida entre $y=x$, el eje $x$ y $0\le x\le2$ tiene área
\[A=\int_0^2 x\dd=2.\]
Sus momentos para $\rho=1$ son
\[M_y=\int_0^2x^2\dd=\frac83,\qquad
M_x=\frac12\int_0^2x^2\dd=\frac43.\]
Por tanto,
\[\boxed{(\bar x,\bar y)=\left(\frac43,\frac23\right)}.\]
El centroide corresponde al promedio de los tres vértices $(0,0)$, $(2,0)$ y $(2,2)$.
\end{ejemplo}
\begin{ejemplo}[Centroide entre dos curvas]
Considere la región entre $y=x$ y $y=x^2$ para $0\le x\le1$. Su área ya fue calculada: $A=1/6$. Los momentos con densidad unitaria son
\begin{align*}
M_y&=\int_0^1 x(x-x^2)\dd
=\left[\frac{x^3}{3}-\frac{x^4}{4}\right]_0^1=\frac1{12},\\
M_x&=\frac12\int_0^1(x^2-x^4)\dd
=\frac12\left(\frac13-\frac15\right)=\frac1{15}.
\end{align*}
Así,
\[\boxed{\bar x=\frac{1/12}{1/6}=\frac12,\qquad
\bar y=\frac{1/15}{1/6}=\frac25}.\]
La ordenada $\bar y$ queda entre las dos curvas, como exige la interpretación geométrica.
\end{ejemplo}

%=========================================================

\subsection*{Por qué los momentos deben dividirse entre la masa}
Para una lámina homogénea entre $g(x)$ y $f(x)$, una tira vertical de área aproximada $(f-g)\Delta x$ tiene centro aproximadamente en la altura $[f(x)+g(x)]/2$. Por eso su contribución al momento respecto del eje $x$ es
\[
\Delta M_x\approx\rho\frac{f(x)+g(x)}2[f(x)-g(x)]\Delta x
=\frac\rho2(f(x)^2-g(x)^2)\Delta x.
\]
Para el momento respecto del eje $y$, la distancia al eje es aproximadamente $x$, de modo que $\Delta M_y\approx\rho x(f-g)\Delta x$. Al integrar y dividir por la masa $m=\rho A$ se obtienen las fórmulas anteriores. Esta deducción también explica que $M_x$ contiene \emph{cuadrados} de las alturas.
\begin{ejemplo}[Centroide de un cuarto de disco]
La región del primer cuadrante limitada por $x^2+y^2\le R^2$ tiene frontera superior $f(x)=\sqrt{R^2-x^2}$ e inferior $g(x)=0$, con $0\le x\le R$. Por simetría sabemos que $\bar x=\bar y$. Su área es $A=\pi R^2/4$. El momento respecto del eje $y$ (con densidad unitaria) es
\[
M_y=\int_0^R x\sqrt{R^2-x^2}\,dx
=\left[-\frac13(R^2-x^2)^{3/2}\right]_0^R=\frac{R^3}{3}.
\]
De donde
\[
\boxed{\bar x=\bar y=\frac{4R}{3\pi}}.
\]
Ambas coordenadas son positivas y menores que $R$, en concordancia con la geometría.
\end{ejemplo}

%=========================================================


\section{Integrales impropias}
%=========================================================
La integral de Riemann estudiada inicialmente se refiere a funciones acotadas en intervalos cerrados. Una integral se denomina \emph{impropia} cuando el intervalo es infinito o el integrando no está acotado cerca de algún extremo o punto interior. En estos casos no basta aplicar una primitiva: hay que \emph{formular y evaluar el límite} correspondiente.
\subsection*{Intervalos infinitos}
\begin{definicion}
Si $f$ es integrable en todo intervalo $[a,b]$ con $b>a$,
\[\boxed{\int_a^{\infty} f(x)\dd:=\lim_{b\to\infty}\int_a^b f(x)\dd,}\]
si el límite existe y es finito. Si no es finito, la integral diverge. De manera semejante se define $\int_{-\infty}^{b}f(x)\dd$.
\end{definicion}
\begin{ejemplo}[Convergencia en un intervalo infinito]
Evaluemos
\[I=\int_1^\infty\frac{1}{x^2}\dd.
\]
Por definición,
\[I=\lim_{b\to\infty}\int_1^b x^{-2}\dd
=\lim_{b\to\infty}\left[-\frac1x\right]_1^b
=\lim_{b\to\infty}\left(1-\frac1b\right)=\boxed{1}.\]
La región bajo $1/x^2$ se extiende sin fin, pero su área total es finita.
\end{ejemplo}
\begin{ejemplo}[Divergencia de la integral armónica]
Para $\int_1^\infty\frac1x\dd$, 
\[\lim_{b\to\infty}\int_1^b\frac1x\dd
=\lim_{b\to\infty}\ln b=+\infty.\]
La integral \emph{diverge}; escribir simplemente $\ln\infty$ no constituye una evaluación matemática.
\end{ejemplo}
\begin{teorema}[Integral $p$ en el infinito]
Para $p\in\R$,
\[\boxed{\int_1^\infty\frac{1}{x^p}\dd\text{ converge si y sólo si }p>1.}\]
\end{teorema}
\begin{proof}
Si $p\ne1$, para $b>1$,
\[\int_1^b x^{-p}\dd=\frac{b^{1-p}-1}{1-p}.\]
Si $p>1$, entonces $1-p<0$ y $b^{1-p}\to0$, con valor límite $1/(p-1)$. Si $p<1$, el término $b^{1-p}$ crece sin cota. Para $p=1$ ya se observó que $\ln b\to\infty$.
\end{proof}

\subsection*{Integrandos no acotados}
Cuando la función diverge cerca de $a$ pero es continua para $a<x\le b$, definimos
\[\int_a^b f(x)\dd:=\lim_{t\downarrow a}\int_t^b f(x)\dd,\]
si el límite es finito. Para una singularidad en el interior $c\in(a,b)$ se deben evaluar \emph{separadamente} las dos integrales laterales.
\begin{ejemplo}[Una singularidad integrable]
\[\int_0^1\frac1{\sqrt{x}}\dd
=\lim_{t\downarrow0}\left[2\sqrt x\right]_t^1
=\lim_{t\downarrow0}(2-2\sqrt t)=\boxed{2}.\]
Aunque $1/\sqrt x$ crece sin cota en $0$, su integral converge.
\end{ejemplo}
\begin{ejemplo}[La integral impropia no admite cancelación artificial]
Examinemos $\int_{-1}^1\frac1x\dd$. La función tiene una singularidad interior en $x=0$. Se requiere que converjan tanto $\int_{-1}^0\frac1x\dd$ como $\int_0^1\frac1x\dd$; sin embargo ambas divergen (a $-\infty$ y $+\infty$). Por tanto, la integral impropia \textbf{diverge}.

El valor principal de Cauchy $\lim_{\varepsilon\downarrow0}\bigl[\int_{-1}^{-\varepsilon}\!dx/x+\int_{\varepsilon}^1\!dx/x\bigr]=0$ es un concepto distinto y no convierte la integral impropia ordinaria en convergente.
\end{ejemplo}
\begin{teorema}[Integral $p$ cerca del origen]
Para $p\in\R$,
\[\boxed{\int_0^1\frac{1}{x^p}\dd\text{ converge si y sólo si }p<1.}\]
\end{teorema}
\begin{proof}
Para $p\ne1$, $\int_t^1x^{-p}\dd=[x^{1-p}/(1-p)]_t^1=(1-t^{1-p})/(1-p)$. Al tomar $t\downarrow0$, este límite es finito precisamente si $1-p>0$. Para $p=1$ surge $-\ln t\to\infty$.
\end{proof}
\begin{ejemplo}[Comparación de integrales impropias]
En $x\ge1$ tenemos $0<\frac1{x^2+1}\le\frac1{x^2}$. Como $\int_1^\infty x^{-2}\dd=1$, el criterio de comparación para integrales no negativas prueba que
\[\int_1^\infty\frac{1}{x^2+1}\dd\]
converge, incluso sin calcular primero su primitiva. Como comprobación adicional, por $\arctan x$ su valor es $\pi/4$.
\end{ejemplo}
\begin{observacion}
La convergencia es una pregunta independiente de la posibilidad de encontrar una antiderivada elemental. Una integral impropia puede converger aunque su primitiva no se exprese mediante funciones elementales.
\end{observacion}

%=========================================================

\subsection*{Criterio de comparación: no siempre hay que encontrar una primitiva}
Para $0\le f(x)\le g(x)$ en $[a,\infty)$, si $\int_a^\infty g(x)\,dx$ converge, también converge $\int_a^\infty f(x)\,dx$, pues sus integrales truncadas son crecientes y están acotadas. De manera contraria, si $0\le g(x)\le f(x)$ y la integral de $g$ diverge, también lo hace la de $f$.
\begin{ejemplo}[Una integral convergente sin primitiva elemental]
Considere
\[
I=\int_0^\infty e^{-x^2}\,dx.
\]
En $[0,1]$ el integrando es continuo y, por ello, integrable. Para $x\ge1$ se cumple $x^2\ge x$ y así $0<e^{-x^2}\le e^{-x}$. Como $\int_1^\infty e^{-x}\,dx=e^{-1}<\infty$, el criterio de comparación garantiza que $I$ converge. Este argumento no requiere hallar una primitiva elemental de $e^{-x^2}$.
\end{ejemplo}
\begin{observacion}
En una integral impropia con una singularidad en un punto interior, las dos integrales laterales deben converger por separado. El valor principal de Cauchy puede existir aunque la integral impropia ordinaria diverja: son conceptos diferentes.
\end{observacion}

%=========================================================


\section{Modelación y verificación con R}
%=========================================================
El software no sustituye el planteamiento. Primero se definen región, variable, límites, radios y unidades; después se verifica la integral. Utilizamos funciones de \texttt{R} base, sin bibliotecas auxiliares.
\begin{ejemplo}[Verificación de un área y de un volumen]
Para el área entre $y=x$ y $y=x^2$ y el volumen de cascarones al girar $y=x^2$ sobre $0\le x\le2$ alrededor del eje $y$, el código es:
\par\medskip
\begin{verbatim}
# Area entre dos graficas
f <- function(x) x - x^2
integrate(f, lower = 0, upper = 1)

# Volumen por cascarones
h <- function(x) 2*pi*x*(x^2)
integrate(h, lower = 0, upper = 2)

# Grafica de la region plana
x <- seq(0, 1, length.out = 200)
plot(x, x, type = "l", xlab = "x", ylab = "y")
lines(x, x^2, lty = 2)
legend("topleft", legend = c("y=x", "y=x^2"), lty = 1:2)
\end{verbatim}
Los resultados deben aproximarse a $1/6$ y $8\pi$, respectivamente. Una discrepancia importante suele indicar un error de planteamiento, no una deficiencia de la herramienta de integración.
\end{ejemplo}
\begin{ejemplo}[Una comprobación de longitud por cuadratura]
Para $y=(2/3)x^{3/2}$ en $[0,3]$:
\par\medskip
\begin{verbatim}
fprima <- function(x) sqrt(x)
L <- integrate(function(x) sqrt(1 + fprima(x)^2), 0, 3)
print(L$value)
print(14/3)
\end{verbatim}
La aproximación numérica coincide con el valor exacto $14/3$.
\end{ejemplo}

%=========================================================

\begin{ejemplo}[Una regla de punto medio programada desde cero]
Esta función sencilla de \texttt{R} sirve para aproximar la integral de una función sin recurrir directamente a \texttt{integrate}; ilustra el paso desde sumas finitas hasta integrales:
\begin{verbatim}
punto_medio <- function(f, a, b, n) {
  h <- (b - a)/n
  x <- a + ((1:n) - 0.5)*h
  h*sum(f(x))
}
f <- function(x) x^2
for (n in c(2, 4, 8, 16, 32)) {
  print(c(n = n, aprox = punto_medio(f, 0, 2, n)))
}
print(8/3)  # Integral exacta de x^2 entre 0 y 2
\end{verbatim}
El valor exacto es $8/3$. Al aumentar $n$ la regla del punto medio se aproxima a ese número. La herramienta verifica el modelo; no determina por sí sola qué variable, límites o sección se deben utilizar.
\end{ejemplo}

%=========================================================


\section{Banco de ejercicios}
%=========================================================
Los reactivos están agrupados por técnica y ordenados de menor a mayor dificultad. Conviene dibujar cada región o sólido antes de integrar. En los problemas físicos, declare unidades y especifique qué representa cada diferencial. Los nuevos reactivos están redactados para estas notas, con una progresión de ejercicios cortos, planteamientos geométricos y problemas de modelación. La metodología se inspira en la construcción por sumas de Riemann de Stewart y en una exposición algebraica rigurosa de los resultados.

\begin{Ejer}\textbf{Área bajo curvas y cambio de signo.} Calcule el área geométrica, no sólo la integral con signo.
\begin{enumerate}[series=banco3]
\item Determine el área entre $y=2x$ y el eje $x$ para $0\le x\le3$.
\item Determine el área entre $y=x^2$ y el eje $x$ para $-1\le x\le2$.
\item Determine el área entre $y=x-2$ y el eje $x$ para $0\le x\le4$.
\item Determine el área entre $y=3-x$ y el eje $x$ para $1\le x\le5$.
\item Determine el área entre $y=x^2-1$ y el eje $x$ para $-2\le x\le2$.
\item Determine el área entre $y=\sin x$ y el eje $x$ para $0\le x\le2\pi$.
\item Determine el área entre $y=x(x-2)$ y el eje $x$ para $0\le x\le3$.
\item Compare $\int_{-1}^1x^3\dd$ con el área entre $y=x^3$ y el eje $x$.
\item Calcule $\int_{-2}^2(x+1)\,dx$ y el área geométrica entre $y=x+1$ y el eje $x$ sobre ese intervalo; compare ambos resultados.
\item Determine el área bajo $y=|x-2|$ en $[0,4]$ mediante integrales sin valor absoluto.
\item Trace la región limitada por $y=2-\lvert x\rvert$ y el eje $x$ en $[-2,2]$ y encuentre su área.
\item Explique geométricamente por qué el área bajo $y=x$ en $[-1,1]$ es positiva aunque la integral definida sea cero.
\end{enumerate}
\end{Ejer}

\begin{Ejer}\textbf{Áreas entre dos gráficas.} Halle todas las intersecciones y determine cuál curva está arriba.
\begin{enumerate}[resume=banco3]
\item $y=2x$ y $y=x^2$.
\item $y=4-x^2$ y $y=0$.
\item $y=x+2$ y $y=x^2$.
\item $y=2-x^2$ y $y=x$.
\item $y=x^3$ y $y=x$ sobre $[-1,1]$.
\item $y=\sqrt x$ y $y=x$ sobre $[0,1]$.
\item $y=\sin x$ y $y=\cos x$ en $[0,\pi/2]$.
\item $y=x^2-2$ y $y=-x$ sobre sus intersecciones.
\item $y=|x|$ y $y=x^2$ entre sus puntos de corte.
\item $y=x^2-4$ y $y=0$ sobre $[-3,3]$.
\item Halle el área entre $y=\cos x$ y $y=\sin x$ en $[0,\pi/2]$, separando en el punto de intersección.
\item Halle el área entre $y=3x-x^2$ y $y=x$ sobre sus intersecciones.
\item Calcule el área entre $y=\lvert x\rvert$ y $y=1$ sobre $[-1,1]$.
\item Para $f(x)=x^2$ y $g(x)=2-x$, justifique cuál curva está arriba en cada subintervalo de $[-2,2]$ y plantee el área total.
\end{enumerate}
\end{Ejer}

\begin{Ejer}\textbf{Integración respecto de $y$.} Exprese el área como $\int(\text{derecha}-\text{izquierda})\,dy$.
\begin{enumerate}[resume=banco3]
\item $x=y^2$ y $x=2y$, $0\le y\le2$.
\item $x=4-y^2$ y $x=0$, $-2\le y\le2$.
\item $x=1+y$ y $x=y^2+1$ sobre sus intersecciones.
\item $x=y^2$ y $x=y+2$ sobre sus intersecciones.
\item $x=\sqrt y$ y $x=y$, $0\le y\le1$.
\item $x=3$ y $x=y^2$, $-\sqrt3\le y\le\sqrt3$.
\item $x=2-y$ y $x=0$, $0\le y\le2$.
\item $x=1-y^2$ y $x=y^2-1$ sobre sus intersecciones.
\item Calcule el área delimitada por $x=y^2$ y $x=4$, integrando respecto de $y$.
\item Para la región encerrada por $x=y^2$ y $x=y+2$, encuentre primero las intersecciones y después el área.
\item Calcule el área entre $x=2y-y^2$ y $x=0$ para $0\le y\le2$.
\item Compare el planteamiento con $dx$ y con $dy$ para la región entre $x=y^2$ y $x=y$, $0\le y\le1$.
\end{enumerate}
\end{Ejer}


\begin{Ejer}\textbf{Volúmenes por secciones transversales no circulares.} Dibuje la base, obtenga la longitud que determina cada sección y formule $A(x)$.
\begin{enumerate}[resume=banco3]
\item La base es $0\le y\le x$, $0\le x\le2$; las secciones perpendiculares al eje $x$ son cuadrados.
\item La base es $0\le y\le\sqrt x$, $0\le x\le4$; las secciones son cuadrados.
\item La base es $0\le y\le x$, $0\le x\le2$; las secciones son triángulos equiláteros.
\item La base es $0\le y\le1+x$, $0\le x\le1$; las secciones son triángulos rectángulos isósceles cuyo cateto es $1+x$.
\item La base es $0\le y\le\sqrt{4-x^2}$, $-2\le x\le2$; cada sección es un semicírculo cuyo diámetro es la longitud vertical de la base.
\item La base es el disco $x^2+y^2\le4$; las secciones son cuadrados cuyo lado es la cuerda vertical de la base.
\item La base es el disco $x^2+y^2\le9$; las secciones son semicírculos cuyo diámetro es la cuerda vertical correspondiente.
\item Explique por qué el volumen por secciones es $\int A(x)\,dx$ incluso cuando el sólido no tiene eje de revolución.
\item La base es la región entre $y=x^2$ y $y=2x$, $0\le x\le2$; las secciones son cuadrados.
\item En el reactivo anterior, sustituya las secciones cuadradas por triángulos equiláteros y compare los volúmenes.
\item La base es $0\le x\le1$, $0\le y\le x^2$; cada sección transversal es un rectángulo cuya altura es el doble de la longitud vertical de la base.
\item Dibuje un sólido con base plana rectangular y secciones transversales de área variable, y escriba una integral consistente con sus dimensiones.
\end{enumerate}
\end{Ejer}



\begin{Ejer}\textbf{Volúmenes de revolución: discos.} Dibuje el eje de giro, identifique el radio y deduzca la integral del volumen.
\begin{enumerate}[resume=banco3]
\item Gire la región bajo $y=x$ sobre $[0,2]$ alrededor del eje $x$.
\item Gire la región bajo $y=x^2$ sobre $[0,1]$ alrededor del eje $x$.
\item Gire la región bajo $y=\sqrt x$ sobre $[0,4]$ alrededor del eje $x$.
\item Gire la región bajo $y=2-x$ sobre $[0,2]$ alrededor del eje $x$.
\item Gire la región comprendida entre $x=y^2$ y el eje $y$, para $0\le y\le2$, alrededor del eje $y$.
\item Gire la región comprendida entre $x=\sqrt y$ y el eje $y$, para $0\le y\le4$, alrededor del eje $y$.
\item Gire la región bajo $y=\sin x$, $0\le x\le\pi$, alrededor del eje $x$.
\item Gire la región bajo $y=1+x$, $0\le x\le2$, alrededor del eje $x$.
\item La base de un sólido es $0\le y\le x$, $0\le x\le2$, y las secciones perpendiculares al eje $x$ son discos cuyo diámetro es $x$. Halle el volumen.
\item La región bajo $y=\sqrt{4-x^2}$, $0\le x\le2$, gira alrededor del eje $x$. Calcule el volumen de la semiesfera resultante.
\item Halle el volumen generado al girar la región entre $y=2$ y el eje $x$, $0\le x\le3$, alrededor del eje $x$.
\item Compare el volumen del sólido obtenido por giro de $y=x$, $0\le x\le1$, con el del cilindro circunscrito de radio $1$ y altura $1$.
\end{enumerate}
\end{Ejer}

\begin{Ejer}\textbf{Volúmenes de revolución: arandelas y ejes desplazados.} Indique explícitamente $R$ y $r$.
\begin{enumerate}[resume=banco3]
\item Gire la región entre $y=2$ y $y=x$, $0\le x\le2$, alrededor del eje $x$.
\item Gire la región entre $y=3$ y $y=x^2$, $0\le x\le1$, alrededor del eje $x$.
\item Gire la región entre $y=x$ y $y=x^2$, $0\le x\le1$, alrededor del eje $x$.
\item Gire la región entre $y=x$ y $y=x^2$, $0\le x\le1$, alrededor de $y=-1$.
\item Gire la región bajo $y=x$, $0\le x\le2$, alrededor de $y=3$.
\item Gire la región entre $x=3$ y $x=y$, $0\le y\le2$, alrededor del eje $y$.
\item Gire la región entre $x=y+1$ y $x=1$, $0\le y\le2$, alrededor del eje $y$.
\item Gire la región entre $y=2x$ y $y=x^2$, $0\le x\le2$, alrededor del eje $x$.
\item La región entre $y=x$ y $y=x^2$, $0\le x\le1$, gira alrededor de $y=2$. Calcule $R$, $r$ y el volumen.
\item La región entre $y=1$ y $y=2$, $0\le x\le3$, gira alrededor del eje $x$. Halle el volumen.
\item Un rectángulo $0\le x\le2$, $-1\le y\le1$, gira alrededor del eje $x$. Explique por qué los cortes son discos y calcule el volumen.
\item La región entre $x=y$ y $x=y^2$, $0\le y\le1$, gira alrededor de la recta $x=2$. Calcule el volumen por arandelas.
\end{enumerate}
\end{Ejer}

\begin{Ejer}\textbf{Cascarones cilíndricos.} Especifique radio, altura y espesor de cada cascarón.
\begin{enumerate}[resume=banco3]
\item Gire la región bajo $y=x$, $0\le x\le2$, alrededor del eje $y$.
\item Gire la región bajo $y=x^2$, $0\le x\le2$, alrededor del eje $y$.
\item Gire la región bajo $y=4-x^2$, $0\le x\le2$, alrededor del eje $y$.
\item Gire la región entre $y=x$ y $y=x^2$, $0\le x\le1$, alrededor del eje $y$.
\item Gire la región bajo $y=2x$, $0\le x\le1$, alrededor de la recta $x=2$.
\item Gire la región bajo $y=3-x$, $0\le x\le3$, alrededor de la recta $x=4$.
\item Gire la región limitada por $x=y^2$, $x=0$, $0\le y\le2$, alrededor del eje $x$, usando cascarones horizontales.
\item Resuelva el volumen del reactivo anterior por discos e indique por qué ambos resultados son iguales.
\item Calcule por cascarones el volumen generado al girar alrededor del eje $y$ la región bajo $y=\sqrt x$ en $0\le x\le4$.
\item Calcule por cascarones horizontales el volumen generado al girar la región $y^2\le x\le y$, $0\le y\le1$, alrededor del eje $x$.
\item Gire la región $0\le y\le x^2$ en $[0,1]$ alrededor de $x=2$. Determine radio, altura y volumen.
\item Para la región entre $y=x$ y $y=x^2$ en $[0,1]$, calcule el volumen al girar alrededor del eje $y$ usando cascarones y verifique con arandelas.
\end{enumerate}
\end{Ejer}

\begin{Ejer}\textbf{Longitud de arco.} Comience encontrando la derivada antes de sustituir en la fórmula.
\begin{enumerate}[resume=banco3]
\item $y=x$, $0\le x\le3$.
\item $y=3x-1$, $0\le x\le2$.
\item $y=\frac23x^{3/2}$, $0\le x\le3$.
\item $y=\frac13(1+x^2)^{3/2}$, $0\le x\le1$.
\item $y=\ln(\cos x)$, $0\le x\le\pi/4$.
\item $x=\frac23y^{3/2}$, $0\le y\le3$.
\item Halle la longitud de un cuarto de circunferencia de radio $2$ por integración.
\item Plantee, sin exigir primitiva elemental, la longitud de $y=x^2$ en $[0,1]$ y aproxime con \texttt{integrate} de R.
\item Halle la longitud de $y=\tfrac{x^2}{2}-\tfrac14\ln x$ para $1\le x\le2$.
\item Plantee y aproxime la longitud de $y=\sin x$ entre $0$ y $\pi/2$; explique si es necesaria una primitiva elemental.
\item Determine la longitud de la curva dada por $x=\tfrac12y^2$ entre $y=0$ y $y=1$.
\item Razone por qué $\int_0^1\sqrt{1+[f'(x)]^2}\,dx\ge1$ para toda $f$ derivable con derivada continua.
\end{enumerate}
\end{Ejer}

\begin{Ejer}\textbf{Superficie de revolución.} Distinga la fórmula de superficie de la de volumen.
\begin{enumerate}[resume=banco3]
\item Gire $y=x$, $0\le x\le2$, alrededor del eje $x$.
\item Gire $y=2$, $0\le x\le3$, alrededor del eje $x$.
\item Gire $y=\sqrt x$, $1\le x\le4$, alrededor del eje $x$.
\item Gire $y=x$, $0\le x\le2$, alrededor del eje $y$.
\item Gire $y=3x$, $0\le x\le1$, alrededor del eje $x$.
\item Gire $y=1+x$, $0\le x\le2$, alrededor del eje $x$.
\item Gire el segmento $y=\sqrt{4-x^2}$, $0\le x\le2$, alrededor del eje $x$ e interprete la superficie.
\item Formule el área de superficie obtenida al girar $y=x^2$ en $[0,1]$ alrededor del eje $x$ y evalúe numéricamente.
\item Calcule el área lateral de la semiesfera obtenida al girar $y=\sqrt{4-x^2}$, $0\le x\le2$, alrededor del eje $x$.
\item Compare superficie lateral y superficie total del cilindro de radio $2$ y altura $3$.
\item Calcule el área de superficie al girar $y=x$ en $[0,1]$ alrededor del eje $y$.
\item Plantee el área de superficie generada al girar $y=x^2+1$, $0\le x\le1$, alrededor del eje $x$, e indique por qué no debe utilizarse $\pi\int f^2$.
\end{enumerate}
\end{Ejer}

\begin{Ejer}\textbf{Trabajo.} Construya el trabajo elemental $dW$ y explique las unidades.
\begin{enumerate}[resume=banco3]
\item Un resorte con $k=80\,\mathrm{N/m}$ se estira $0.10\,\mathrm m$ desde su longitud natural. Calcule el trabajo.
\item Un resorte con $k=100\,\mathrm{N/m}$ se estira desde $x=0.10$ hasta $x=0.30\,\mathrm m$. Calcule el trabajo.
\item Se aplica $F(x)=3x^2+2$ newtons entre $x=0$ y $x=2\,\mathrm m$. Calcule el trabajo.
\item Se aplica $F(x)=10-x$ newtons entre $x=0$ y $x=4\,\mathrm m$. Calcule el trabajo.
\item Un cable pesa $8\,\mathrm{N/m}$ y tiene $5\,\mathrm m$ de longitud. Halle el trabajo para elevarlo completamente hasta su extremo superior.
\item Una cadena de $3\,\mathrm m$ pesa $12\,\mathrm{N/m}$. Halle el trabajo necesario para subirla completamente.
\item Un depósito de sección constante $2\,\mathrm{m}^2$ está lleno de agua hasta $3\,\mathrm m$. Calcule el trabajo para bombearla al borde (use $\rho=1000\,\mathrm{kg/m^3}$ y $g=9.8\,\mathrm{m/s^2}$).
\item Un depósito de profundidad $2\,\mathrm m$ tiene sección horizontal $A(y)=1+y\,\mathrm{m}^2$ a la altura $y$ medida desde el fondo. Formule el trabajo para bombear el agua al borde.
\item La fuerza resistente es $F(x)=-4x\,\mathrm N$ en $[0,3]\,\mathrm m$. Halle el trabajo de la resistencia y el de un agente externo que compense esa fuerza.
\item Un tanque tiene sección horizontal $A(y)=2y\,\mathrm{m^2}$, $0\le y\le2$, y está lleno de agua. Halle el trabajo para bombear el agua hasta $y=2$.
\item Para un gas idealizado con presión $P(V)=k/V$, $V_1<V_2$, encuentre el trabajo de expansión y sus unidades.
\item Un cable uniforme de $6\,\mathrm m$ pesa $5\,\mathrm{N/m}$. Represente con una integral el trabajo para levantar únicamente los $3\,\mathrm m$ inferiores hasta la parte superior.
\end{enumerate}
\end{Ejer}


\begin{Ejer}\textbf{Valor promedio y teorema del valor medio para integrales.} Calcule el promedio y, cuando proceda, halle $c$ tal que $f(c)=f_{\mathrm{prom}}$.
\begin{enumerate}[resume=banco3]
\item $f(x)=x$, $0\le x\le2$.
\item $f(x)=x^2$, $0\le x\le2$.
\item $f(x)=3x+1$, $0\le x\le4$.
\item $f(x)=2-x^2$, $-1\le x\le1$.
\item $f(x)=\sin x$, $0\le x\le\pi$.
\item $f(x)=e^x$, $0\le x\le1$.
\item $f(x)=1/x$, $1\le x\le e$.
\item Una partícula tiene velocidad $v(t)=t^2+1$ $\mathrm{m/s}$, $0\le t\le3$ segundos. Halle el desplazamiento y la velocidad promedio.
\item Una temperatura varía según $T(t)=18+4t-t^2$ $^\circ\mathrm C$ durante $0\le t\le4$ horas. Calcule el promedio.
\item Justifique por qué la altura $f_{\mathrm{prom}}$ de un rectángulo de base $b-a$ produce la misma área con signo que la integral de $f$.
\item Dé un ejemplo de función continua que toma su valor promedio en más de un punto del intervalo.
\item Utilice la regla del punto medio con $n=4$ para estimar el promedio de $f(x)=x^2$ en $[0,2]$ y compárelo con el valor exacto.
\end{enumerate}
\end{Ejer}



\begin{Ejer}\textbf{Centro de masa.} Identifique masa, momentos y coordenadas del centroide.
\begin{enumerate}[resume=banco3]
\item Una varilla uniforme ocupa $[0,4]$ y tiene densidad $3$. Encuentre su masa y centro.
\item Una varilla ocupa $[0,2]$ y tiene densidad $\rho(x)=1+x$. Encuentre $m$ y $\bar x$.
\item Una varilla ocupa $[0,1]$ y tiene densidad $\rho(x)=2x+1$. Encuentre $m$ y $\bar x$.
\item Una varilla ocupa $[1,3]$ y tiene densidad $\rho(x)=x$. Encuentre $m$ y $\bar x$.
\item Encuentre el centroide del triángulo bajo $y=x$, $0\le x\le2$.
\item Encuentre el centroide de la región bajo $y=x^2$, $0\le x\le1$.
\item Encuentre el centroide de la región entre $y=x$ y $y=x^2$, $0\le x\le1$.
\item Encuentre el centroide de la región bajo $y=2-x$, $0\le x\le2$.
\item Calcule el centroide de un cuarto de disco de radio $R$ en el primer cuadrante.
\item Calcule el centroide bajo $y=\sqrt x$, $0\le x\le1$, con densidad superficial uniforme.
\item Una varilla ocupa $[0,3]$ y tiene densidad $\rho(x)=2+x$. Obtenga masa, momento y centro.
\item Justifique que para una lámina uniforme entre dos funciones $f\ge g$ la fórmula de $M_x$ contiene la diferencia $f^2-g^2$.
\end{enumerate}
\end{Ejer}

\begin{Ejer}\textbf{Integrales impropias y criterios de convergencia.} Escriba el límite que define la integral antes de evaluarla.
\begin{enumerate}[resume=banco3]
\item $\displaystyle\int_1^\infty x^{-2}\dd$.
\item $\displaystyle\int_1^\infty x^{-1}\dd$.
\item $\displaystyle\int_1^\infty x^{-3/2}\dd$.
\item $\displaystyle\int_2^\infty e^{-x}\dd$.
\item $\displaystyle\int_0^1x^{-1/2}\dd$.
\item $\displaystyle\int_0^1x^{-3/2}\dd$.
\item $\displaystyle\int_0^1\ln x\dd$.
\item $\displaystyle\int_{-1}^1\frac1x\dd$.
\item $\displaystyle\int_1^\infty\frac{1}{x^2+1}\dd$.
\item $\displaystyle\int_0^\infty e^{-2x}\dd$.
\item Determine para qué $p$ converge $\displaystyle\int_2^\infty x^{-p}\dd$.
\item Determine para qué $q$ converge $\displaystyle\int_0^2x^{-q}\dd$.
\item Determine si converge $\displaystyle\int_1^\infty\frac1{x^2+x}\,dx$ y calcule su valor.
\item Demuestre por comparación que $\displaystyle\int_1^\infty e^{-x^2}\,dx$ converge.
\item Determine para qué valores de $p$ converge $\displaystyle\int_0^1x^{-p}\,dx$ y compare con el criterio en el infinito.
\item Estudie la convergencia de $\displaystyle\int_{-1}^2\frac1{x-1}\,dx$ separando correctamente los límites laterales.
\end{enumerate}
\end{Ejer}

\begin{Ejer}\textbf{Integradores y selección del modelo.} En cada caso entregue planteamiento, solución e interpretación.
\begin{enumerate}[resume=banco3]
\item Una región entre $y=x$ y $y=x^2$ en $[0,1]$ gira alrededor del eje $x$. Calcule el volumen por arandelas y luego por cascarones horizontales.
\item La misma región gira alrededor del eje $y$. Resuelva por cascarones y compruebe por arandelas horizontales.
\item Una fuerza $F(x)=4x$ newtons se opone al desplazamiento de $x=0$ a $x=3$. Halle el trabajo realizado \emph{por la fuerza} y el trabajo que debe suministrar un agente externo para movimiento cuasiestático.
\item Una región bajo $y=2-x$, en $0\le x\le2$, gira alrededor del eje $x$. Halle volumen y superficie del sólido, y explique la diferencia física.
\item Calcule la longitud de $y=x$ entre $0$ y $1$, y la superficie que produce al girar alrededor del eje $x$.
\item Sea $f(x)=x^{-2}$ en $[1,\infty)$. Compare el área bajo la gráfica con el volumen obtenido al girar alrededor del eje $x$ (demuestre la convergencia de ambos).
\item Esboce el modelo de trabajo para extraer agua de un tanque cuya sección horizontal es $A(y)=y+1$ entre $0\le y\le2$ y verifique las unidades del integrando.
\item Explique por qué $\int_{-1}^1dx/x$ no es cero como integral impropia, aunque la función sea impar.
\item Una base plana es $0\le x\le2$, $0\le y\le x$; sus secciones perpendiculares al eje $x$ son cuadrados. Deduzca y calcule el volumen.
\item Los valores de $v_A-v_B$ medidos en $t=1,3,5,7$ segundos son $1,3,2,1$ $\mathrm{m/s}$. Estime la separación acumulada de dos vehículos durante $[0,8]$ utilizando la regla del punto medio.
\item Una temperatura varía como $T(t)=20+2t$ grados Celsius, $0\le t\le4$ horas. Halle la temperatura promedio y el instante en que se alcanza.
\item Diseñe un ejemplo en que el cambio de variable de $dx$ a $dy$ evite dividir la región en dos integrales; justifique el criterio geométrico.
\end{enumerate}
\end{Ejer}

\clearpage
\section*{Síntesis para el planteamiento}
\begin{formula}[Qué se acumula en cada aplicación]
\begin{align*}
\text{Área entre curvas: }& A=\int(\text{altura de tira})\,dx,\\
\text{Volumen por secciones: }&V=\int(\text{área de sección})\,dx,\\
\text{Cascarones: }&V=2\pi\int(\text{radio})(\text{altura})\,dx,\\
\text{Longitud: }&L=\int\sqrt{1+[f'(x)]^2}\,dx,\\
\text{Superficie: }&S=2\pi\int(\text{radio})\,ds,\\
\text{Trabajo: }&W=\int(\text{fuerza})\,dx,\\
\text{Momento: }&M=\int(\text{posición})\,dm,\\
\text{Valor promedio: }&f_{\mathrm{prom}}=\frac{1}{b-a}\int_a^b f(x)\,dx.
\end{align*}
Estas expresiones son un recordatorio de la \emph{construcción}. Los límites, signos, radios y dimensiones se determinan en cada problema antes de integrarlos.
\end{formula}

\end{document}
